\documentclass[11pt]{article}
\usepackage{mzqastyle}
\setrunningtitle{AI Analyst — qlib Integration}

\begin{document}

\mzqatitlepage
  {Engineering integration reference}
  {qlib Integration}
  {Return \& risk prediction, portfolio optimization, and score enhancement}
  {AI Analyst embeds Microsoft \code{qlib} as a lean, in-process library to power
   its quantitative desk: a machine-learned return forecast, a factor-structured
   forward covariance, portfolio optimization, and a walk-forward backtest, with
   the learned signal folded into the idea scanner and the investment committee.
   This volume documents the \emph{engineering} of that integration --- the
   lean-embed principle, the module map, the Postgres-to-qlib panel adapter, the
   dual optimizer dispatcher, the committee and scanner wiring, the REST surface,
   and operations. The \emph{mathematics} of each model is specified in the
   \emph{Quantitative \& ML Models} volume.}
  {qlib Integration \quad·\quad Generated 31 July 2026}

\mzqatoc
\clearpage

%======================================================================
\section{What the integration delivers}

The quantitative desk adds a machine-learned layer on top of the deterministic
valuation engine and the fundamentals warehouse:

\begin{itemize}
  \item \textbf{Return prediction} --- a cross-sectional LightGBM alpha model over
        the monthly fundamentals panel, producing a per-name expected forward
        return (the \code{mu} consumed by the optimizer and the scanner).
  \item \textbf{Risk prediction} --- a factor-structured forward covariance
        (\code{StructuredCovEstimator}), $\Sigma=F\,\Sigma_B\,F\tr+\diag(\mathbf d)$,
        with its decomposition exposed to the enhanced-indexing optimizer.
  \item \textbf{Portfolio optimization} --- both a native optimizer \emph{and}
        qlib's optimizers, selectable at runtime.
  \item \textbf{Score enhancement} --- the learned alpha plus regime-conditioned
        factor-IC weights blended into the value $+$ sentiment interest score.
  \item \textbf{Committee integration} --- a \code{quant\_signals} evidence node and
        a Quantitative Factor Analyst specialist.
  \item \textbf{REST $+$ UI} --- a \code{/api/quant/*} router and a \emph{Quant} tab.
\end{itemize}

Everything degrades gracefully: with no trained model, or with qlib or an optional
dependency absent, each surface returns a defined fallback and the application runs
unchanged.

\begin{keybox}[The lean-embed principle]
qlib's model, risk estimator, optimizers, and evaluation run as \emph{plain
libraries} on a MultiIndex DataFrame the app builds from Postgres. The integration
never uses qlib's expression engine, its \code{.bin} data provider, the Alpha158/360
handlers, or the recorder/\code{qrun} workflow --- so there is no \code{qlib.init},
no data-provider bootstrap, and no mlflow tracking backend. Only \code{lightgbm},
\code{scikit-learn}, and \code{cvxpy} become hard new runtime dependencies.
\end{keybox}

\section{How qlib maps onto the desk's needs}

\begin{center}\small
\begin{tabular}{@{}>{\bfseries\color{navy}}l >{\ttfamily\footnotesize}p{5.4cm} >{\ttfamily\footnotesize}l@{}}
\toprule
\tabhead{Need} & \normalfont\tabhead{qlib capability used} & \normalfont\tabhead{Module} \\
\midrule
Return prediction      & LGBModel $+$ DatasetH $+$ DataHandlerLP.from\_df & api/quant/qlib\_alpha.py \\
Risk prediction        & StructuredCovEstimator & api/quant/qlib\_risk.py \\
Portfolio optimization & PortfolioOptimizer, EnhancedIndexingOptimizer & api/quant/qlib\_optimize.py \\
Evaluation / backtest  & calc\_ic, risk\_analysis & api/quant/qlib\_backtest.py \\
\bottomrule
\end{tabular}
\end{center}

\section{Module map}

\begin{center}\small
\begin{tabular}{@{}>{\ttfamily\footnotesize\color{navy}}l p{9.5cm}@{}}
\toprule
\normalfont\tabhead{File} & \normalfont\tabhead{Responsibility} \\
\midrule
qlib\_data.py     & Build the monthly cross-sectional feature/label panel from Postgres (reuses the regime-IC loaders and point-in-time alignment). \\
qlib\_alpha.py    & \code{AlphaLGB} model (recorder-free), \code{train}/\code{predict}/\code{evaluate}, persistence, lazy cache, CLI. \\
qlib\_risk.py     & \code{structured\_cov()} $\to$ annualized $\Sigma$ $+$ the $(F,\Sigma_B,\mathbf d)$ decomposition; sync price/returns loader. \\
qlib\_optimize.py & \code{solve(optimizer, \dots)} dispatcher over the native $+$ five qlib backends $\to$ a uniform \code{PortfolioSolution}. \\
qlib\_backtest.py & Walk-forward out-of-sample signal backtest scored with qlib \code{risk\_analysis}. \\
alpha\_signal.py  & Central bridge: cached latest cross-section $+$ per-ticker expected returns (used by scanner, committee, router). \\
ic\_weights.py    & Regime-conditioned factor-IC $\to$ per-family scoring weights. \\
optimize.py / risk.py & Native SLSQP/MIP optimizer; sample / Ledoit--Wolf / FF structured covariance. \\
routers/quant.py  & \code{/api/quant/\{backends,alpha,risk,optimize,backtest\}}. \\
\bottomrule
\end{tabular}
\end{center}

\begin{figure}[h]
\centering
\begin{tikzpicture}[node distance=7mm and 10mm, font=\sffamily\footnotesize]
  \node[store, minimum width=28mm] (pg) {Postgres\\\scriptsize \code{fact\_metrics\_*}};
  \node[proc, right=of pg, minimum width=24mm] (panel) {qlib\_data\\\scriptsize build\_panel (PIT, $z$-scored)};
  \node[llm, right=of panel, minimum width=27mm, align=left] (models)
    {\textbf{qlib\_alpha} $\to\ \muv$\\[1pt]\textbf{qlib\_risk} $\to\ \Sigma$};
  \node[llm, right=of models, minimum width=24mm] (opt) {qlib\_optimize\\\scriptsize native $|$ qlib $\to\ \wv$};
  \node[data, right=of opt, minimum width=20mm] (sol) {Portfolio\\solution};

  \node[proc, below=11mm of panel, minimum width=24mm] (bridge) {alpha\_signal\\\scriptsize cached cross-section};
  \node[proc, below=11mm of models, minimum width=32mm, xshift=4mm] (scan) {scanner $+$ committee\\\scriptsize interest score · quant node};

  \draw[arr] (pg) -- (panel);
  \draw[arr] (panel) -- (models);
  \draw[arr] (models) -- (opt);
  \draw[arr] (opt) -- (sol);
  \draw[arr] (panel) -- (bridge);
  \draw[arr] (bridge) -- (scan);
  \draw[arr] (models.south) to[out=-90,in=90] (scan.north);
\end{tikzpicture}
\caption{The lean qlib embed. One point-in-time panel feeds the alpha and risk
models; the optimizer dispatcher consumes their $\muv$ and $\Sigma$; a cached bridge
exposes the alpha to the scanner and committee.}
\end{figure}

\section{The panel adapter}
\code{qlib\_data.build\_panel} produces a qlib-style MultiIndex frame --- rows
$(\text{datetime},\text{instrument})$, columns
$(\text{feature}\,|\,\text{label})$ --- consumed directly by
\code{DataHandlerLP.from\_df}. Features are importance-ranked fundamentals from four
families, cross-sectionally $z$-scored per month; the label is the forward 1m/3m
return; fundamentals are aligned point-in-time ($\text{period\_end}+90\text{d}\to$
month-end). This is the \emph{same} loader and alignment the regime-IC research
uses, so the production forecast and the research signal share one lookahead-safe
data path. The alpha model trains through the real qlib \code{DatasetH} pipeline;
the \code{AlphaLGB} subclass attaches LightGBM's callbacks directly in \code{fit},
so training runs with plain LightGBM and needs no \code{qlib.init} or tracking
backend. Prediction runs the stored booster on any reindexed feature matrix.

\section{Dual optimizer dispatch}
\code{qlib\_optimize.solve(optimizer, tickers, mu, sigma, \dots)} dispatches to
either backend family and returns a uniform \code{PortfolioSolution} (weights,
annualized expected return / vol / Sharpe, factor exposures, diagnostics):

\begin{center}\small
\begin{tabular}{@{}>{\ttfamily\footnotesize\color{navy}}l p{5.0cm} p{5.0cm}@{}}
\toprule
\normalfont\tabhead{\ttfamily optimizer} & \normalfont\tabhead{Engine} & \normalfont\tabhead{Notes} \\
\midrule
native & \code{optimize.py} (SLSQP / cvxpy-MIP) & Full constraints: sector caps, vol cap, cardinality, turnover. \\
qlib\_mvo / \_gmv / \_rp / \_inv & qlib \code{PortfolioOptimizer} & Mean-variance / min-variance / risk-parity / inverse-vol; long-only, fully invested. \\
qlib\_enhanced\_indexing & qlib \code{EnhancedIndexingOptimizer} & Benchmark-relative tracking-error; consumes the \code{StructuredCovEstimator} decomposition. \\
\bottomrule
\end{tabular}
\end{center}

Both families are fed the same annualized $\muv$ (alpha or historical mean) and
$\Sigma$ (qlib-structured, Ledoit--Wolf, or sample). Neither is a default-only path
--- both ship enabled and are chosen per request or from the UI selector.

\section{Score enhancement: alpha $+$ regime-IC weights}
The value $+$ sentiment interest score (\code{screener\_agent}) blends two learned
inputs on top of the fundamental composite:

\begin{enumerate}
  \item \textbf{Normalized alpha} $v_\alpha$ --- the qlib model's expected return,
        cross-sectionally normalized, as the dominant term (default weight $0.40$)
        when the model covers the shortlist.
  \item \textbf{Regime-IC family weights} --- the fundamental score budget
        (value / growth) is re-split by how predictive each family has been
        \emph{in the current macro regime} (\S\ref{sec:regime}), rather than by
        fixed constants.
\end{enumerate}

When no trained model or IC data is available the score reduces exactly to the
fixed-weight composite (guaranteed by test). Each row reports \code{alpha},
\code{alpha\_percentile}, and a \code{score\_components} breakdown; the response
carries a \code{scoring} block naming the weights and model actually used.

\section{Where the macro-regime IC feeds forecasting}
\label{sec:regime}
The macro-cycle model (HMM / VAE) labels each month with a regime; the regime-IC
engine then computes, \emph{conditioned on that regime}, the Spearman rank IC of
every fundamental metric against forward returns, into
\code{fact\_equity\_factor\_ic\_regime}. The one live consumption path in the desk
is scoring: \code{ic\_weights.family\_weights} reads the latest regime's rows,
aggregates per-metric $|\text{IC}|$ to family-level emphasis, and the scanner
re-splits its fundamental budget accordingly --- so the interest score leans harder
on whichever factor family the current regime says is most predictive.

\begin{notebox}[Scope]
The LightGBM alpha model itself is regime-\emph{agnostic}: it learns from factor
\emph{values} (including momentum/volatility from the market-factor family), and
the regime enters forecasting through the score weights, not as a model feature.
The committee receives macro regime context from the macro node and the
\code{quant\_signals} block, but not the factor-IC table directly.
\end{notebox}

\section{Committee integration}
\begin{itemize}
  \item \textbf{State} --- \code{InvestmentCommitteeState.quant\_signals} plus a
        config toggle \code{enable\_quant\_signals}.
  \item \textbf{Evidence node} --- \code{qlib\_signals\_node} runs in the
        \emph{prepare} phase (once, shared across providers): the ticker's alpha
        expected return, universe percentile, and model rank-IC, plus a peer-group
        factor-risk read (forward vol, factor exposures, minimum-variance weight).
        It auto-fans-out from the engine and auto-fans-in to every analyst, and
        degrades to \code{available: false} on any error, so it never fails a run.
  \item \textbf{Specialist} --- the Quantitative Factor Analyst reconciles the
        statistical alpha/risk signals with the fundamental/DCF thesis and flags
        alpha-vs-valuation disagreements.
\end{itemize}

\section{REST API and frontend}
Router \code{api/routers/quant.py} (under \code{/api/quant}), all handlers
offloading qlib work to worker threads:

\begin{center}\small
\begin{tabular}{@{}>{\ttfamily\color{navy}}l p{9.4cm}@{}}
\toprule
\normalfont\tabhead{Endpoint} & \normalfont\tabhead{Purpose} \\
\midrule
GET /backends   & Optimizer backends, risk models, alpha-model metadata. \\
POST /alpha     & Expected returns for a universe (or the latest top-$N$ cross-section). \\
POST /risk      & Forward covariance / volatility / factor exposures. \\
POST /optimize  & Portfolio optimization; \code{optimizer} selects native or any qlib backend. \\
POST /backtest  & Walk-forward out-of-sample signal backtest with qlib \code{risk\_analysis}. \\
\bottomrule
\end{tabular}
\end{center}

The \emph{Quant} tab (\code{QuantView.tsx}) exposes an optimizer-backend selector,
risk-model and $\mu$-source selectors, a live expected-returns table, an
optimized-weights table, and a backtest panel.

\section{Operations}

\begin{itemize}
  \item \textbf{qlib on Python 3.13.} No PyPI wheel exists for 3.13, so qlib is
        built from a bundled clone and installed non-editable, so \code{import qlib}
        resolves from any working directory (the C++ toolchain is required for the
        build). Lean runtime deps are pinned in \code{backend/requirements.txt}
        (\code{lightgbm}, \code{scikit-learn}, \code{cvxpy}, \code{ecos},
        \code{dill}, \code{pyarrow}, \dots).
  \item \textbf{Cold start / pre-warm.} The first alpha call for a jurisdiction
        builds the cross-section panel from the warehouse and then caches it (TTL
        via \code{QLIB\_ALPHA\_CACHE\_TTL}); the walk-forward backtest is cached
        separately. Both are pre-warmed on startup so the first user request is
        served from cache.
  \item \textbf{Retraining.} The alpha model is (re)trained via the CLI:
        \code{python -m api.quant.qlib\_alpha train --jurisdiction US}; artifacts
        are persisted per jurisdiction and hot-reloaded on file mtime.
  \item \textbf{Tests.} \code{backend/api/quant/tests/test\_quant.py} covers the
        optimizer backends, the structured covariance, the exact scoring fallback,
        the alpha blend, the IC re-split, and a DB-gated panel/train smoke test.
\end{itemize}

\vfill
\begin{center}\small\color{muted}
A lean, in-process qlib embed: one point-in-time panel, a learned alpha and a
structured covariance, dual optimization, and a walk-forward backtest --- wired into
the scanner and the committee, and degrading cleanly to the prior behavior wherever
a model or dependency is absent.
\end{center}

\end{document}
